% Solutions to the Chapter 8 exercises that are not code, from lectures/L08/appendix/c_exercises.md.
% Exercises 8.6 and 8.7 are Code and are not answered; see chapter.tex for why.

\section{Chapter 8: Interrupts and Deferred Work}
\label{sol:sec:c8}
Answers to the exercises in \secref{c8:sec:exercises}. Exercises~\ref{c8:ex:6} and~\ref{c8:ex:7}
are Code, and are checked on the target as their Check yourself lines say.

% ------------------------------------------------------------------------------------------
\solution[fillaccent2]{c8:ex:1}{The three numbers}{Recall}

\xpart{a} The property has three cells. \code{0} says the interrupt is an SPI, a shared peripheral
interrupt; \code{112} is its number among the SPIs; \code{4} is the trigger. On a real board whoever
wrote the device tree copied 112 out of the SoC's manual. Here QEMU writes it: it places
\code{qa-dev} on the virt machine's platform bus, whose interrupts start at SPI 112, and generates
the node with the number it assigned. \textbf{144} is the same interrupt as the GIC numbers it, the
identity the controller signals to the CPU, and the hardware column of \file{/proc/interrupts}.
\textbf{17} is the Linux IRQ number, handed out when the GIC's IRQ domain mapped hardware interrupt
144 for the driver: an index into the kernel's own table, derived from neither of the others
(\secref{c8:app:a1}).

\xpart{b} Where the SPIs begin in the GIC's own numbering. Interrupts 0 to 15 are SGIs, the
software-generated interrupts CPUs send one another, and 16 to 31 are PPIs, private to each CPU,
the per-CPU timer among them. The device tree counts within the SPIs and the GIC counts all of
them, so the GIC driver adds 32 when it translates an SPI's cell, and 16 for a PPI's
(\file{drivers/irqchip/irq-gic.c}).

\xpart{c} \textbf{The Linux number.} It is handed out in the order mappings are created, so an
extra device whose interrupt is mapped earlier in the boot takes a number and moves everything
after it. The DT cell and the GIC number describe the wiring, which an extra device does not
change. One qualification for this machine: QEMU hands out the platform bus's interrupts to dynamic
devices in the order they are created, first free first, so an extra device of that kind created
before \code{qa-dev} on the command line would move 112 and 144 as well.

\xpart{d} Two mistakes. 112 is the device tree's number and appears nowhere in
\file{/proc/interrupts}, which shows the Linux number, 17, first and the GIC's, 144, in its hardware
column; the thing to grep for is the name passed to \code{request_irq}, \code{qa_irq}. And a
missing line does not mean a missing driver: a line appears only while a handler is registered for
the interrupt, so a loaded module that has not requested it, or has already freed it, is absent
too. Whether the module is loaded is a question for \code{lsmod}.

\xpart{e} \code{4} is \code{IRQ_TYPE_LEVEL_HIGH}: level-sensitive, active high, so the interrupt is
pending for as long as the line is high. \code{1} would have been \code{IRQ_TYPE_EDGE_RISING}: the
controller notices the low-to-high transition and nothing else. For a GIC SPI these are the only
two allowed. Configured as an edge, \code{qa-dev}, which holds its line high until its status is
cleared, would lose any event that arrived while the line was still high, because no new edge
happens (\secref{c8:app:a4}).

% ------------------------------------------------------------------------------------------
\solution[fillaccent2]{c8:ex:2}{Reading \code{/proc/interrupts}}{Recall}

\xpart{a} \code{17}, the Linux IRQ number. \code{1789}, how many times this interrupt was taken on
CPU0, and \code{0}, on CPU1. \code{GIC-0}, the interrupt controller, the first GIC. \code{144}, the
hardware number, the GIC's own. \code{Level}, the trigger type. \code{qa_irq}, the name passed to
\code{request_irq}; there is one per registered handler, separated by commas on a shared line
(\secref{c8:app:a2}).

\xpart{b} \textbf{Nothing is wrong.} The GIC delivers an SPI to one CPU at a time, and Linux points
it at the first online CPU in the interrupt's affinity mask. The mask allows both CPUs by default,
so CPU0 gets every one, and nothing on this minimal target, no \code{irqbalance}, ever moves it.
\textbf{Or something has taken CPU1 out of it}: the mask was narrowed to CPU0, by hand or with the
\code{irqaffinity=} boot parameter, or CPU1 is not taking interrupts at all. To tell them apart,
read \file{/proc/irq/17/smp_affinity_list}, which CPUs are allowed, and
\file{/proc/irq/17/effective_affinity_list}, where the interrupt actually goes: \code{0-1} and
\code{0} is the first explanation. CPU1's column for the other interrupts, the timer and the IPIs,
shows whether it takes interrupts at all. Then write \code{1} to \file{smp_affinity_list} and watch
CPU1's count for interrupt 17 start to climb.

\xpart{c} Ruled out: anything in your handler, which is not running, and a missing
acknowledgement, which would make the count race rather than stand still. Still possible:
everything between the device and the CPU. The driver asked for the wrong interrupt, from the wrong
node or the wrong index in it; the trigger type is wrong; the line is masked in the controller; or
the device is not asserting at all, because \code{ENABLE} or \code{IRQ_EN} is not set, or because
\code{PERIOD_NS} is below \code{QA_DEV_MIN_PERIOD_NS}, 1,000 ns, at which the device refuses to
start. And a count that stands still need not be zero: it may be left over from an earlier load.

\xpart{d} The source of the interrupt is never cleared. On a level-triggered line the handler
returns, the line is still asserted, and the handler is entered again at once, as fast as the CPU
can take it: the missing acknowledgement of \secref{c8:app:a6}. On \code{qa-dev} that is a handler
that neither writes \code{IRQ_STATUS} nor empties the FIFO.

\xpart{e} No. The counts belong to the interrupt descriptor, not to your registration, and
\code{free_irq} does not reset them: the run in \secref{c8:app:a2} read 919 after the first load
and 993 immediately after the second. So any rate has to be a difference between two readings,
divided by the time measured between them; a single reading counts everything since the descriptor
was created, across every load.

% ------------------------------------------------------------------------------------------
\solution[fillaccent3]{c8:ex:3}{The acknowledgement}{Hand calculation}

\xpart{a} Every 1 written clears its bit, so \code{SAMPLE} is cleared, which was wanted. The
\code{OVERRUN} that arrived a microsecond after your read is cleared too, and the same write resets
the sticky \code{OVERRUN} in \code{STATUS}: the device dropped a sample and nothing in the driver
will ever know. The other thirty bits do not exist, and the ones written to them do nothing.

\xpart{b} \code{SAMPLE} is cleared exactly as before. The late \code{OVERRUN} is not: its bit was 0
in the value written, and writing 0 leaves a bit alone. So it stays set, the line stays asserted,
since the device raises it while either bit is set, and the handler is entered again at once,
reads \code{0b10}, and deals with the overrun. Nothing is lost.

\xpart{c} With a bit that only the write clears:
\begin{enumerate}
  \item The handler reads \code{0b01}.
  \item It writes \code{0b01} back. The bit clears and the line drops.
  \item A new sample arrives: it is queued, its bit is set again, and the line rises.
  \item The handler drains the FIFO, taking the new sample with the others.
  \item It returns \code{IRQ_HANDLED}, with the bit still set and the FIFO empty.
  \item The line is still asserted, so the handler runs again: it reads \code{0b01}, acknowledges,
    finds nothing to drain, and returns.
\end{enumerate}
One interrupt with nothing to do, and \textbf{no data lost}: the sample was taken in step 4. On
\code{qa-dev} itself even that interrupt does not happen. The read that empties the FIFO clears
\code{SAMPLE} (part \textbf{f)}), so step 5 returns with the bit clear and the line low, and the
controller counts a level-triggered interrupt as pending only while its line is high, so the rise
in step 3 leaves nothing behind.

\xpart{d}
\begin{enumerate}
  \item The handler reads \code{0b01}.
  \item It drains the FIFO.
  \item A new sample arrives after the last read of the FIFO: it is queued, its bit is set, and the
    line rises.
  \item The handler writes the \code{0b01} it read in step 1, which clears the bit the new sample
    set, and the line drops.
  \item It returns \code{IRQ_HANDLED}, and nothing is pending.
\end{enumerate}
The sample is in the FIFO and nothing will interrupt for it. It is collected only when some later
event sets the bit again: on \code{qa-dev}, the next sample, one period late; on a device whose
events are sporadic, whenever the next one comes, if one does.

\xpart{e} \textbf{c) is survivable and d) is not.} c) costs work: one extra entry to the handler
that finds nothing, which is visible, limits itself, and loses no sample. d) destroys information:
the only notice of an event is cleared, its data waits on some unrelated later event, and nothing
anywhere reports it. That is why \secref{c8:app:a5} acknowledges exactly what it read, and does it
before the drain.

\xpart{f} \textbf{Empty the FIFO.} Reading \code{QA_DEV_SAMPLE} pops one sample, and the pop that
takes the level to zero clears \code{IRQ_SAMPLE} and lowers the line, without the driver writing
\code{IRQ_STATUS} at all. To find it, search \file{tools/qa-dev/qa-dev.c} for every assignment to
\code{irq_status}. It is set in \code{qa_dev_timer_expired}; cleared by a write to
\code{IRQ_STATUS} in \code{qa_dev_write}; zeroed by \code{CTRL_RESET}, which throws the samples
away with it; and cleared in \code{qa_dev_fifo_pop} when \code{fifo_level} reaches zero, which is
the answer. It covers \code{SAMPLE} only: \code{OVERRUN} is cleared by a write to
\code{IRQ_STATUS}, or a reset, and by nothing else.

% ------------------------------------------------------------------------------------------
\solution[fillaccent3]{c8:ex:4}{An interrupt budget}{Hand calculation}

\xpart{a} One every 250 us is $1{,}000{,}000 / 250 = 4{,}000$ interrupts a second. At 8 us each,
$4{,}000 \times 8 = 32{,}000$ us, 32 ms of every second: \textbf{3.2\%} of one CPU.

\xpart{b} The rate is unchanged, 4,000 a second, and $4{,}000 \times 40 = 160{,}000$ us, 160 ms of
every second: \textbf{16\%}.

\xpart{c} \textbf{250 us}, where $4{,}000 \times 250$ us is the whole second. In practice sooner,
because taking and leaving an interrupt costs time the handler's own figure leaves out, and nothing
else runs on that CPU meanwhile. Past 250 us the handler cannot finish one interrupt before the next
is due, and the backlog grows until the device's FIFO overflows.

\xpart{d} Sixteen samples at one per 250 us fill the FIFO in $16 \times 250 = 4{,}000$ us:
\textbf{4 ms}. Counting from the last drain, the worst case has the first new sample arrive at once,
the sixteenth fill the FIFO 3.75 ms later, and the seventeenth, at 4 ms, dropped. On this course's
kernel that is exactly one jiffy.

\xpart{e} The hard handler is now 10 us: $4{,}000 \times 10 = 40$ ms a second, \textbf{4\%} of a
CPU with interrupts disabled. The threaded half is $4{,}000 \times 30 = 120$ ms a second, 12\%, in a
thread the scheduler can place, preempt and prioritise. What does not change is the rate, 4,000 a
second; the point at which the CPU saturates, because each interrupt still costs 40 us of a 250 us
period; and the 4 ms of part \textbf{d)}, which belongs to the device's FIFO and rate, although if
the thread does the draining it is now the thread's scheduling delay that must stay under it.
\textbf{The system's total work does not fall.} It rises a little: still 40 us an interrupt, 16\% of
a CPU, plus waking the thread and switching to it and back, on every interrupt. What moved is who
waits: 12\% of a CPU went from time nothing can preempt to time the scheduler decides about
(\secref{c8:app:b6}).

% ------------------------------------------------------------------------------------------
\solution[fillaccent4]{c8:ex:5}{Where does the work go}{Design}

\xpart{a} \textbf{Hard handler.} It is a couple of register accesses, only the handler can do it,
and until it is done the line stays asserted.

\xpart{b} \textbf{Threaded handler}, with \code{IRQF_ONESHOT} so that the line stays masked while it
runs. 4 KB through a register is 1,024 reads of 32 bits, long enough to matter (\secref{c8:app:b5}),
and a FIFO that deep is itself the buffer that lets the copy wait for the thread. The case for the
hard handler is the one in \secref{c8:app:b1}, work that would be lost if it waited, and it has to
show that the thread's scheduling delay could outlast the FIFO.

\xpart{c} \textbf{Whichever half puts the data in the buffer, straight after putting it there.}
\code{wake_up_interruptible} does not sleep and is safe in a hard handler (\secref{c9:app:a4}), so
in the lab's design, where the hard handler drains, the hard handler. Waking before the data is
there wakes a reader to an empty buffer.

\xpart{d} \textbf{Threaded handler}, or a workqueue if the allocation is not tied to one interrupt:
\code{GFP_KERNEL} may sleep while the kernel reclaims memory. Better still, allocate in advance.

\xpart{e} \textbf{Workqueue}, as delayed work queued 10 ms ahead. The retry is not one per
interrupt, and a threaded handler that slept for 10 ms would keep the line masked, and the next
interrupt waiting, for all of it.

\xpart{f} \textbf{Hard handler}: an \code{atomic_inc}, or an increment under the lock the handler
already holds, costs nanoseconds.

\xpart{g} \textbf{Threaded handler.} \code{dev_err} is legal in a hard handler, because
\code{printk} never sleeps, but it is not cheap: the message may be written out to a serial console
before the call returns, at 115,200 baud about 87 us a character. A hard handler that can log on
every interrupt should at least use \code{dev_err_ratelimited}.

\xpart{h} \textbf{d)} and \textbf{e)}. Both reach \code{schedule()} from a context that has no
task of its own to put to sleep. A \code{GFP_KERNEL} allocation is marked as a place that may
sleep, and under memory pressure it does, entering reclaim; \code{msleep} always does. In a hard
handler the scheduler finds itself called with interrupts disabled in the middle of an interrupt,
reports \code{BUG: scheduling while atomic}, and the machine is left with an interrupt that never
completes. The allocation is the worse of the two, because it sleeps only when memory is short, so
the bug waits for a machine under load.

\xpart{i} \code{CONFIG_DEBUG_ATOMIC_SLEEP}, which \file{kernel/qa.config} enables. It makes every
call that may sleep check its context, and prints \code{BUG: sleeping function called from invalid
context} with a backtrace at the call, every time, including the allocations that would not have
slept that time.

% ------------------------------------------------------------------------------------------
\solution[fillaccent4]{c8:ex:8}{The interrupt that never stops}{Design}

\xpart{a} \code{insmod} never returns. The storm starts inside the module's init, the moment it
sets \code{IRQ_EN}, on CPU0, where the interrupt is routed, and the task that was interrupted there,
\code{insmod} in the course's run, never runs again. The console says nothing for over twenty
seconds; in the course's run the first line after the module's own two was the RCU stall report, at
29.8 s of uptime (\secref{c8:app:a6}). The wait is the stall detector's timeout,
\code{CONFIG_RCU_CPU_STALL_TIMEOUT=21} on this kernel, counted from the start of the grace period
that CPU0 is holding up.

\xpart{b} \code{rcu: INFO: rcu_preempt detected stalls on CPUs/tasks:}, then \code{0-....: (0 ticks
this GP)} and \code{Sending NMI from CPU 1 to CPUs 0:}. A grace period ends only when every CPU has
passed through a quiescent state: a context switch, the idle loop, or user mode. CPU0 is in hard
interrupt context, re-entering the handler for ever, and passes through none of them; ``0 ticks
this GP'' says it has not even taken a timer tick. After 21 seconds RCU reports the CPU that is
holding it up, and CPU1, which is still running, asks CPU0 for a backtrace, which is what would name
the handler. RCU is not the cause but the first victim with a timeout; this kernel builds neither
the soft- nor the hard-lockup detector, which would otherwise have spoken first.

\xpart{c} \textbf{TIMEOUT.} \file{ci/test.sh} kills QEMU after 120 seconds and reports ``L08:
TIMEOUT'' with the path to the serial log. A failure means the script ran to the end, printed
\code{QA-TEST-END} with its status, and some checks failed. Here \file{run.sh} never got past
\code{insmod}: there is no status, and nothing is known about any check after it. Calling that a
failure would claim the checks ran; the distinction sends you to the serial log, where the evidence
is, rather than to a list of check results that does not exist.

\xpart{d} \code{grep qa_irq /proc/interrupts}, twice, a second apart, on CPU1, the CPU still
running. It would show interrupt 17's count on CPU0 growing by far more than the thousand a second
the period allows, while CPU1's stays at 0: the racing count of \secref{c8:app:a2}, and the
diagnosis in one line.

\xpart{e} \textbf{No.} An edge-triggered line interrupts on a transition, not a level. \code{qa-dev}
raises its line and holds it until the bit is cleared, so with nothing clearing it there is one
rising edge, one interrupt, and then silence: the line stays high and no new edge ever comes. The
machine is fine and the device looks dead. The FIFO fills in sixteen periods, \code{OVERRUN} sets,
and every later sample is dropped, silently.

% ------------------------------------------------------------------------------------------
\solution[fillaccent]{c8:ex:9}{Interrupts counted two ways}{Cross-check}

\xpart{a} $2\text{ s} / 1\text{ ms} = \mathbf{2{,}000}$.

\xpart{c}
\begin{itemize}
  \item \textbf{A delta}, because the count belongs to the descriptor and survives \code{free_irq}:
    across a reload it carried on from 919 to 993 (\secref{c8:app:a2}). A single reading counts
    everything since boot, and has no interval to divide by.
  \item \textbf{Using the measured time makes the shortfall larger, not smaller.} Take the course's
    hard-handler run, 1,789 interrupts in 2.11 s (\secref{c8:app:b4}). Against 2.00 s the ideal is
    2,000 and the shortfall $211/2{,}000 = 10.6\%$; against 2.11 s the ideal is 2,110 and the
    shortfall $321/2{,}110 = 15.2\%$. It moves up by 4.7 points, because the interval was longer
    than two seconds, so more interrupts were due.
  \item The hypothesis is yours to write down, and part \textbf{d)} is where it meets the source.
\end{itemize}

\xpart{d}
\begin{itemize}
  \item \textbf{Relative.} The line is \code{timer_mod(s->timer,
    qemu_clock_get_ns(QEMU_CLOCK_VIRTUAL) + s->period_ns)}, run from the timer's own callback after
    the sample is queued: the next expiry is ``now plus a period'', where now is when QEMU got round
    to running the callback, not the moment it was due.
  \item So each interval is the period plus the callback's lateness, plus the few microseconds the
    callback spends before it reads the clock again. Every late callback pushes all the later ones
    back, so lateness accumulates rather than averaging out. Note what is not in it: your driver.
    The device re-arms inside its own callback, whether or not the handler has run.
  \item A shortfall $s$ against a period $P$ means each interval is $P/(1-s)$, so the missing term is
    $o = P\,s/(1-s)$. At 15\% and 1,000 us that is $1{,}000 \times 0.15/0.85 = 176$ us; the run above,
    at 15.2\%, gives 179 us, and the Cross-check's own 1 ms row, at 17.6\%, gives 213 us. A couple
    of hundred microseconds is plausible for an emulator whose timers are serviced by threads on a
    busy host, and would be absurd for a hardware timer. Nothing inside the guest can confirm it.
\end{itemize}

\xpart{e}
\begin{itemize}
  \item \textbf{The implied overhead would be the same in every row}, and the shortfall would fall
    as the period grows, as $o/(P+o)$. With the 1 ms row's 213 us: at 10 ms, $213/10{,}213 = 2.1\%$,
    about 206 interrupts in 2.1 s, where 190 were measured; at 0.1 ms, $213/313 = 68\%$, about 6,700
    in 2.1 s, where 13,085 were measured. Too few at the slow end, twice too many at the fast end.
  \item \file{/proc/uptime} is the guest kernel's clock, counted from the generic timer that QEMU
    emulates from the same virtual clock the device's timer runs on. So it contains the kernel,
    whose \code{sleep 2} overran to about 2.1 s, and QEMU. It does not contain your driver, which
    decides neither when interrupts are raised nor how time is kept.
  \item A candidate, not a finding. At 10 ms the guest's CPUs are idle almost all the time and
    QEMU's threads sleep in the host, so each timer waits on a host wake-up, whose lateness is up to
    the host's scheduler. At 0.1 ms the guest is busy, QEMU is awake, and an expired timer is
    noticed sooner. The table's pattern, 947 us of lateness at 10 ms and 63 us at 0.1 ms, fits
    that, and does not prove it.
  \item One to run: read the device's timestamp registers, \code{QA_DEV_TS_LO} and \code{TS_HI}, in
    the handler and difference successive stamps. That is the device's own interval between
    assertions, in the emulator's clock, with the guest's handling taken out; repeat it with the
    other CPU kept busy, and see whether the lateness at 10 ms falls. One you cannot run here: the
    same driver on real hardware with an oscilloscope on the interrupt line, which has no emulator
    in the measurement at all.
\end{itemize}

\xpart{f}
\begin{itemize}
  \item The line is a level. Once \code{SAMPLE} is set it stays high until the handler clears it,
    and a sample that arrives while it is already high changes nothing on the line, so no second
    interrupt is raised; the sample joins the FIFO, and the handler's drain takes everything that
    has accumulated. At 0.1 ms, whenever the time from assertion to drain exceeds 100 us, a second
    sample is waiting: $14{,}389 - 13{,}708 = 681$ extra samples, 1.05 per interrupt.
  \item At 1 ms the handler always drains well inside a period, so every assertion finds exactly
    one sample. The period changed against the time from assertion to drain; the driver did not.
  \item \code{OVERRUN}: bit 1 of \code{IRQ_STATUS}, and the sticky bit 3 of \code{STATUS}. The
    device sets it when a sample arrives to find the FIFO full, which would take sixteen undrained
    samples, 1.6 ms without service at this period. In the course's run it did not happen: the
    interrupts coalesced and no data was lost (\secref{c8:app:b6}). To check a run of your own,
    record the bit in the handler, or run the device in counter mode, where sample $n$ is $n$, and
    look for a gap: a dropped sample still advances the device's count.
\end{itemize}

\xpart{g} Wrong: \textbf{``dropping''}, when nothing is dropped: the device raises fewer interrupts,
because each interval is the period plus the emulator's lateness, and every interrupt it raises is
serviced. \textbf{``15\%''}, when the shortfall depends on the period, 8.7\%, 17.6\% and 38.6\% at
10, 1 and 0.1 ms, and even two runs at 1 ms differ, 15.2\% and 17.6\%. \textbf{``the driver is
broken''}, when the driver's two counters agree with the device and the cause is in the device
model's re-arm and the emulator. Supported by the course's measurements: ``At \code{PERIOD_NS} =
1,000,000, \code{qa-dev} raised 1,789 interrupts in 2.11 s of guest uptime, 15.2\% fewer than the
2,110 the period implies, and the driver serviced every one of them, one sample each. The device
re-arms its timer from when its callback ran, so each interval is the period plus the emulator's
lateness, about 180 us here, and that lateness is not constant across periods.'' Still unknown: why
the lateness depends on the period; how much of it is QEMU, the host, or the guest's load; whether
real hardware shows any of it; and anything about latency at all, since every figure here is a count
or an average.

\xpart{h} Not a number from any of this: counts say how many interrupts were raised, not how long
any of them waited. What belongs on a datasheet is a distribution of the interval from the line's
assertion to the handler, over many thousands of interrupts, with its mean, high percentiles and
the worst value observed, and the conditions beside it: kernel version and configuration
(6.12.30, \code{PREEMPT}, \code{HZ=250}), hard or threaded handler, interrupt rate, system load,
number of samples and duration, and the platform, which here is an emulator. Refuse to claim a
guaranteed worst case, any number for real hardware, and anything inferred from counts or averages.
The register is \code{QA_DEV_TS_LO}/\code{TS_HI}, which records the instant the device asserted
its line, in the emulator's clock that the guest's own clock also runs on. Latency is the interval
from that instant to the handler, and software can read when its handler ran but not when the line
went up; on real hardware only an instrument on the wire sees that. The register hands software the
start of the interval, so it can subtract instead of infer.
